% TODO checklist:
% - [x] Inspected src/observability/tracing.py (OpenTelemetry integration)
% - [x] Inspected src/observability/middleware.py (trace ID propagation)
% - [x] Reviewed Q&A.md sections on tracing and OpenTelemetry
% - [x] Reviewed README.md for tracing overview
% - [x] Examined OTLP exporter configuration options

\section{Tracing}
\label{sec:tracing}

The Cineca Agentic Platform integrates OpenTelemetry for distributed tracing, providing end-to-end visibility into request flows across service boundaries. The tracing subsystem is designed with graceful degradation---applications run normally even when OpenTelemetry dependencies are unavailable or tracing is disabled.

\subsection{OpenTelemetry Integration Architecture}

The platform's tracing implementation wraps the OpenTelemetry SDK with a defensive layer that handles missing dependencies, configuration errors, and runtime failures without impacting application functionality:

\begin{lstlisting}[language=Python, caption={Defensive OpenTelemetry initialization (src/observability/tracing.py)}]
try:
    from opentelemetry import trace as otel_trace
    from opentelemetry.sdk.resources import Resource
    from opentelemetry.sdk.trace import TracerProvider
    from opentelemetry.sdk.trace.export import BatchSpanProcessor
    from opentelemetry.sdk.trace.sampling import (
        AlwaysOnSampler, ParentBased, TraceIdRatioBased
    )
    _OTEL_AVAILABLE = True
except Exception:
    _OTEL_AVAILABLE = False

def setup_tracing(app=None) -> bool:
    """
    Initialize OpenTelemetry tracing if enabled in configuration.
    Returns True if tracing is active, False if no-op.
    """
    global _INITIALIZED, _PROVIDER
    
    if _INITIALIZED:
        return _PROVIDER is not None
    
    if not getattr(settings, "OTEL_ENABLED", False):
        log.info("otel.tracing.disabled")
        _INITIALIZED = True
        return False
    
    if not _OTEL_AVAILABLE:
        log.warning("otel.tracing.requested_but_unavailable")
        _INITIALIZED = True
        return False
    
    # ... proceed with initialization ...
\end{lstlisting}

\subsection{Resource Attributes}

OpenTelemetry traces are enriched with semantic resource attributes that identify the service, version, and deployment context:

\begin{lstlisting}[language=Python, caption={OpenTelemetry resource attribute configuration}]
def _build_resource() -> Resource:
    """Construct a Resource with standard semantic attributes."""
    hostname = socket.gethostname()
    attrs = {
        "service.name": getattr(settings, "APP_NAME", 
                                "cineca-agentic-platform"),
        "service.version": __version__,
        "service.instance.id": hostname,
        "deployment.environment": getattr(settings, "APP_ENV", "dev"),
        "host.name": hostname,
    }
    return Resource.create(attrs)
\end{lstlisting}

These attributes follow OpenTelemetry Semantic Conventions and enable filtering, grouping, and comparison of traces across deployments.

\subsection{Sampling Strategies}

The platform implements environment-aware sampling to balance observability needs with storage and network costs:

\begin{lstlisting}[language=Python, caption={Environment-aware trace sampling}]
def _select_sampler() -> ParentBased:
    """
    Choose a sampler based on environment:
    - Production: ParentBased(TraceIdRatioBased(X)) where X from 
                  OTEL_SAMPLER_RATIO (default 0.2 = 20%)
    - Non-production: ParentBased(AlwaysOnSampler)
    """
    env = getattr(settings, "APP_ENV", "dev").lower()
    if env in {"prod", "production"}:
        ratio = float(getattr(settings, "OTEL_SAMPLER_RATIO", 0.2))
        ratio = max(0.0, min(1.0, ratio))
        return ParentBased(TraceIdRatioBased(ratio))
    return ParentBased(AlwaysOnSampler())
\end{lstlisting}

The sampling configuration:

\begin{table}[ht]
\centering
\caption{Trace sampling configuration by environment.}
\label{tab:trace-sampling}
\begin{tabular}{lll}
\hline
\textbf{Environment} & \textbf{Sampler} & \textbf{Default Ratio} \\
\hline
Development/Test & \texttt{AlwaysOnSampler} & 100\% \\
Production & \texttt{TraceIdRatioBased} & 20\% \\
\hline
\end{tabular}
\end{table}

The \texttt{ParentBased} wrapper ensures that child spans inherit the sampling decision from their parent, maintaining trace continuity across service boundaries.

\subsection{OTLP Exporters}

The platform supports both gRPC and HTTP/protobuf OTLP exporters for sending traces to an OpenTelemetry Collector or compatible backend:

\begin{lstlisting}[language=Python, caption={OTLP exporter configuration}]
def _make_exporter() -> object | None:
    """Create OTLP exporter according to protocol and endpoint."""
    protocol = str(getattr(settings, "OTEL_EXPORTER_OTLP_PROTOCOL", 
                          "grpc")).lower()
    endpoint = str(getattr(settings, "OTEL_EXPORTER_OTLP_ENDPOINT", 
                          "")).strip()
    
    if not endpoint:
        endpoint = ("http://otel-collector:4317" if protocol == "grpc" 
                   else "http://otel-collector:4318/v1/traces")
    
    if protocol in {"grpc", "otlp", "otlp_grpc"}:
        if OTLPGrpcSpanExporter is None:
            log.warning("opentelemetry OTLP gRPC exporter not available")
            return None
        insecure = endpoint.startswith("http://")
        return OTLPGrpcSpanExporter(endpoint=endpoint, insecure=insecure)
    
    # Default to HTTP/protobuf
    if OTLPHttpSpanExporter is None:
        log.warning("opentelemetry OTLP HTTP exporter not available")
        return None
    return OTLPHttpSpanExporter(endpoint=endpoint)
\end{lstlisting}

\begin{table}[ht]
\centering
\caption{OTLP exporter configuration options.}
\label{tab:otlp-config}
\begin{tabular}{lp{8cm}}
\hline
\textbf{Environment Variable} & \textbf{Description} \\
\hline
\texttt{OTEL\_ENABLED} & Enable/disable tracing (default: \texttt{false}) \\
\texttt{OTEL\_EXPORTER\_OTLP\_PROTOCOL} & Protocol: \texttt{grpc} or \texttt{http} (default: \texttt{grpc}) \\
\texttt{OTEL\_EXPORTER\_OTLP\_ENDPOINT} & Collector endpoint (default: \texttt{http://otel-collector:4317}) \\
\texttt{OTEL\_SAMPLER\_RATIO} & Production sampling ratio (default: \texttt{0.2}) \\
\texttt{OTEL\_CONSOLE\_EXPORTER} & Enable console span output for debugging \\
\hline
\end{tabular}
\end{table}

\subsection{Automatic Instrumentation}

When available, the tracing system automatically instruments common libraries and frameworks:

\begin{lstlisting}[language=Python, caption={Automatic instrumentation setup}]
# FastAPI instrumentation
if FastAPIInstrumentor is not None and app is not None:
    try:
        FastAPIInstrumentor().instrument_app(app, tracer_provider=provider)
        log.info("otel.tracing.fastapi.instrumented")
    except Exception:
        log.exception("otel.tracing.fastapi.instrument_failed")

# Requests library instrumentation
if RequestsInstrumentor is not None:
    try:
        RequestsInstrumentor().instrument()
        log.info("otel.tracing.requests.instrumented")
    except Exception:
        log.exception("otel.tracing.requests.instrument_failed")

# Logging instrumentation (inject trace context into logs)
if LoggingInstrumentor is not None:
    try:
        LoggingInstrumentor().instrument(set_logging_format=True)
        log.info("otel.tracing.logging.instrumented")
    except Exception:
        log.exception("otel.tracing.logging.instrument_failed")
\end{lstlisting}

The automatic instrumentation captures:

\begin{itemize}
    \item \textbf{FastAPI}: Incoming HTTP requests with route, method, status code, and exception information.
    \item \textbf{Requests}: Outgoing HTTP calls to LLM providers and external services.
    \item \textbf{Logging}: Trace context injection into log entries for correlation.
\end{itemize}

\subsection{Trace Context Propagation}

The platform propagates trace context through HTTP headers and exposes trace IDs in response headers for client-side correlation:

\begin{lstlisting}[language=Python, caption={Trace ID extraction and response headers}]
def _current_trace_id() -> str | None:
    """Extract current trace ID for response headers."""
    if not otel_trace:
        return None
    try:
        span = otel_trace.get_current_span()
        ctx = span.get_span_context()
        if ctx and getattr(ctx, "trace_id", 0):
            return f"{ctx.trace_id:032x}"
    except Exception:
        return None
    return None

# In middleware:
trace_id = _current_trace_id()
if trace_id:
    bind_contextvars(trace_id=trace_id)
    response.headers["X-Trace-Id"] = trace_id
\end{lstlisting}

Clients can use the \texttt{X-Trace-Id} response header to search for traces in the backend visualization tool (Jaeger, Zipkin, Grafana Tempo).

\subsection{Manual Span Creation}

For custom instrumentation beyond automatic coverage, the platform provides a \texttt{get\_tracer()} function that returns either a real tracer or a no-op implementation:

\begin{lstlisting}[language=Python, caption={Manual tracer access with no-op fallback}]
def get_tracer(name: str) -> otel_trace.Tracer:
    """
    Return a tracer. If tracing is disabled/unavailable, 
    a no-op tracer is returned.
    """
    if not _OTEL_AVAILABLE:
        class _NoopTracer:
            def start_as_current_span(self, *_, **__):
                from contextlib import nullcontext
                return nullcontext()
        return _NoopTracer()
    
    return otel_trace.get_tracer(name, __version__)
\end{lstlisting}

Usage example for custom span creation:

\begin{lstlisting}[language=Python, caption={Example of manual span instrumentation}]
from src.observability.tracing import get_tracer

tracer = get_tracer(__name__)

async def process_agent_run(run_id: str, input_text: str) -> str:
    with tracer.start_as_current_span("agent.run.process") as span:
        span.set_attribute("run.id", run_id)
        span.set_attribute("input.length", len(input_text))
        
        # ... processing logic ...
        
        span.set_attribute("output.tokens", token_count)
        return output
\end{lstlisting}

\subsection{Graceful Shutdown}

The tracing system provides a shutdown function to flush pending spans before process termination:

\begin{lstlisting}[language=Python, caption={Trace provider shutdown}]
def shutdown_tracing() -> None:
    """Flush and shutdown the tracer provider (graceful shutdown)."""
    global _PROVIDER
    if _PROVIDER is not None:
        try:
            _PROVIDER.shutdown()
            log.info("otel.tracing.shutdown.complete")
        except Exception:
            log.exception("otel.tracing.shutdown.failed")
        finally:
            _PROVIDER = None
\end{lstlisting}

This function should be called during application shutdown (e.g., in FastAPI's \texttt{@app.on\_event("shutdown")} handler) to ensure that all buffered spans are exported before the process exits.

\subsection{Trace Visualization and Analysis}

Traces exported via OTLP can be visualized in various backends:

\begin{itemize}
    \item \textbf{Jaeger}: Open-source distributed tracing system with dependency analysis and performance optimization features.
    
    \item \textbf{Grafana Tempo}: Scalable distributed tracing backend integrated with Grafana for unified observability.
    
    \item \textbf{Zipkin}: Distributed tracing system with a focus on latency troubleshooting.
    
    \item \textbf{Cloud Providers}: AWS X-Ray, Google Cloud Trace, and Azure Monitor support OTLP ingestion.
\end{itemize}

\subsection{Tracing Best Practices}

The platform's tracing implementation follows several best practices for effective distributed tracing:

\begin{enumerate}
    \item \textbf{Meaningful Span Names}: Span names use verb.noun patterns (e.g., \texttt{agent.run.execute}, \texttt{llm.call.generate}) that describe the operation being performed.
    
    \item \textbf{Bounded Attribute Values}: Large values (prompts, responses) are summarized or omitted to prevent trace storage bloat; length attributes are recorded instead.
    
    \item \textbf{Error Recording}: Exceptions are recorded on spans with stack traces for debugging, but sensitive error details are filtered.
    
    \item \textbf{Sampling Awareness}: Critical paths (authentication, authorization) may use explicit sampling overrides to ensure they're always captured.
    
    \item \textbf{Trace Context Preservation}: Trace context is propagated across async boundaries and background job execution.
\end{enumerate}

